\documentclass[book.tex]{subfiles}

\begin{document}

\chapter{Multilayer Perceptrons}
\label{chap:mlpd}
\noindent\rule{11cm}{0.4pt} \\
\textbf{Prerequisites:} \autoref{chap:ml_basics},  \\
\textbf{Difficulty Level:} ** \\
\noindent\rule{11cm}{0.4pt}
\vspace{5mm}

The limitation of linear models is expressivity. There are many complex interactions which can't be modeled simply with linear interactions. Practitioners often bypass this limitation by adding in new "features" that encode pairwise interactions, but ideally we would like a model that can natively learn more complex interactions. We want to add some nonlinearity to our function class $\mathcal{G}$. Let $\sigma$ be some nonlinear function. Then we can define a model by

\begin{align} 
h[w,b](x) = \sigma(w^T x + b) \end{align}

There are many choices commonly made for the function $\sigma$. One common choice is the sigmoidal function

\begin{align} 
\sigma(x) = \frac{1}{1+ e^{-x}}
\end{align}

Note that this function ranges from 0 to 1 (see Figure~\ref{fig:logistic}). We can turn this definition into code.

\begin{marginfigure}
    \centering
    % \includegraphics{figures/Differentiable Models/multilayer_perceptrons/logistic_curve.png}
    \includegraphics{figures/Differentiable Models/multilayer_perceptrons/Logisticurve0.png}
    \caption{A plot of the logistic curve}%\textcolor{red}{TODO: Replace.}}
    \label{fig:logistic}
\end{marginfigure}
\begin{lstlisting}[language=python]
def $\sigma$(x: $\mathbb{R}$) $\to$ $\mathbb{R}$:
  1/(1 + exp(-x))

class h(w: $\mathbb{R}^n$, b: $\mathbb{R}$):
  def $\lambda$(x: $\mathbb{R}^n$) $\to$ $\mathbb{R}$:
    $\sigma$(w$^T$ * x + b)
\end{lstlisting}

With this choice of $\sigma$, the function $h_{(w,b)}$ is commonly referred to as a "perceptron." It corresponds to a very crude model of a neuron. The core idea is that this function corresponds to the following "firing" rule

\begin{align} h[b](x) = \begin{cases}
    1 & \sigma(w^T x + b) > 0.5 \\
    0 & \textrm{otherwise} \\
    \end{cases} 
\end{align}
    
Perceptrons spurred a lot of excitement when they were first introduced in the 1950s. Minsky and Papert showed in their book "Perceptrons" \cite{minsky2017perceptrons} that perceptrons were incapable of learning simple functions such as XOR. This caused a broad loss of interest in Perceptrons and neural networks more broadly. Interest only revived when researchers realized that many of the limitations of perceptrons could be overcome by chaining them into multiple layers. 

To implement multilayer perceptrons, we alter our function $h[w, b]$ to return a vector of outputs rather than a scalar. Let $W \in \mathbb{R}^{P \times M}$ and let $c \in \mathbb{R}^{P}$. Then we defined $h[W, c]: \mathbb{R}^M \to \mathbb{R}^P$ as 

\begin{align} h[W, c](x) = W x + c \end{align}

We then introduce a convenient bit of notation. If $x$ is a vector, we let $\sigma(x)$ denote the component wise application of $\sigma$ to every scalar in the vector. We're now ready to introduce a simple multilayer perceptron

\begin{align} h[W_0, c_0, w_1, b_1](x) = \sigma(w_1^T\sigma(W_0 x + c_0) + b_1) \end{align}

\begin{figure}
    \centering
    % \includegraphics{figures/Differentiable Models/multilayer_perceptrons/mlp_example.png}
    \includegraphics{figures/Differentiable Models/multilayer_perceptrons/mlp_example_0.png}
    %\caption{\textcolor{red}{TODO: Replace this with an illustration of our own. We don't own this figure}}
    \label{fig:my_label}
\end{figure}

We say this this model has one \textit{hidden layer} since the inner layer $\sigma(W_0x + c_0)$ isn't directly reflected in the output. We can transform this definition into Physika code in a straightforward fashion.

\begin{lstlisting}
class OneLayerNet($W_0$: $\mathbb{R}$[N], $w_1$: $\mathbb{R}$, $b_0$: $\mathbb{R}$, $b_1$: $\mathbb{R}$):
  def $\lambda$(x: $\mathbb{R}$[N]) $\to$ $\mathbb{R}$:
    $\sigma$($w_1$*$\sigma$($W_0$*x+$b_0$)+$b_1$)
\end{lstlisting}

There's no reason we have to stop at just one layer here. We can go deeper. Here's an example with two hidden layers for example.

\[ h[W_0, c_0, W_1, c_1, w_2, b_2](x) = \sigma(w_1^T\sigma(W_1 \sigma(W_0 x + c_0) + c_1) + b_1)  \]

You might now start to see where the term \textit{deep learning} originates from. As we keep adding more and more hidden layers, we make the model deeper. Let's take a look at some simple Physika code for a $n$ hidden layer network.

\begin{lstlisting}[language=python]
class FullyConnectedNetwork($\sigma$: $\mathbb{R}$ $\to$ $\mathbb{R}$, W: $\mathbb{R}$[n, N, N], B: $\mathbb{R}$[n, N], w: $\mathbb{R}^N$, b: $\mathbb{R}$):
  def $\lambda$(x: $\mathbb{R}$[N]) $\to$ $\mathbb{R}$:
    for i: x = $\sigma$(W[i] * x + B[i])
    w$^T$ x + b
\end{lstlisting}


\end{document}